=\frac{L}{4}.
\]
Así:
\[
\boxed{\bar x=\qty{1}{\metre}}.
\]
\end{solucion}

\begin{ejercicio}[2: Reacciones bajo carga triangular]
Una viga simplemente apoyada de longitud $L=\qty{6}{\metre}$ soporta una carga triangular descendente con intensidad máxima $q_0=\qty{12}{\kilo\newton\per\metre}$ en el apoyo izquierdo y cero en el derecho.

Calcular las reacciones.
\end{ejercicio}

\begin{solucion}
La resultante vale:
\[
Q=\frac12q_0L=\frac12(12)(6)=\qty{36}{\kilo\newton}.
\]
Está situada a:
\[
\bar x=\frac{L}{3}=\qty{2}{\metre}
\]
desde el extremo de máxima intensidad.

Equilibrio vertical:
\[
R_A+R_B=36.
\]
Momentos respecto a $A$:
\[
R_B(6)-36(2)=0.
\]
Luego:
\[
\boxed{R_B=\qty{12}{\kilo\newton}},
\qquad
\boxed{R_A=\qty{24}{\kilo\newton}}.
\]

\begin{tip}
Una comprobación física inmediata: como la carga está concentrada más cerca de $A$, necesariamente debe cumplirse $R_A>R_B$.
\end{tip}
\end{solucion}

\begin{ejercicio}[3: Tensión y deformación nominal/verdadera]
Una probeta tiene:
\[
L_0=\qty{100}{\milli\metre},\qquad A_0=\qty{200}{\milli\metre\squared}.
\]
Durante el ensayo:
\[
L=\qty{112}{\milli\metre},\qquad A=\qty{175}{\milli\metre\squared},\qquad F=\qty{30}{\kilo\newton}.
\]
Calcular tensión nominal, tensión verdadera, deformación nominal y deformación verdadera.
\end{ejercicio}

\begin{solucion}
Tensión nominal:
\[
\sigma_{\mathrm{nom}}=\frac{30000}{200}=\boxed{\qty{150}{\mega\pascal}}.
\]
Tensión verdadera:
\[
\sigma_{\mathrm{real}}=\frac{30000}{175}=\boxed{\qty{171.43}{\mega\pascal}}.
\]
Deformación nominal:
\[
\varepsilon_{\mathrm{nom}}=\frac{112-100}{100}=\boxed{0.12}.
\]
Deformación verdadera:
\[
\varepsilon_{\mathrm{real}}=\ln\left(\frac{112}{100}\right)=\boxed{0.1133}.
\]

\begin{idea}
La diferencia entre ambas medidas comienza a ser apreciable cuando la deformación deja de ser pequeña.
\end{idea}
\end{solucion}

\begin{ejercicio}[4: Barra elástica sometida a axil]
Una barra de aluminio tiene:
\[
A=\qty{600}{\milli\metre\squared},\qquad
L=\qty{1.2}{\metre},\qquad
E=\qty{70}{\giga\pascal}.
\]
Está sometida a $N=\qty{90}{\kilo\newton}$ de tracción.

Calcular $\sigma$, $\varepsilon$ y $\Delta L$.
\end{ejercicio}

\begin{solucion}
Como $1\ \mathrm{N/mm^2}=1\ \mathrm{MPa}$:
\[
\sigma=\frac{90000}{600}=\boxed{\qty{150}{\mega\pascal}}.
\]
Por Hooke:
\[
\varepsilon=\frac{\sigma}{E}=\frac{150}{70000}=\boxed{2.143\times10^{-3}}.
\]
Finalmente:
\[
\Delta L=\varepsilon L.
\]
Trabajando en milímetros:
\[
\Delta L=2.143\times10^{-3}(1200)=\boxed{\qty{2.57}{\milli\metre}}.
\]
\end{solucion}

\begin{ejercicio}[5: Tensor de pequeñas deformaciones]
El campo de desplazamientos plano es:
\[
u=2\times10^{-3}x+1\times10^{-3}y,
\]
\[
v=-0.5\times10^{-3}x+1\times10^{-3}y.
\]
Calcular $\varepsilon_{xx}$, $\varepsilon_{yy}$, $\varepsilon_{xy}$ y $\gamma_{xy}$.
